\chapter{放射监测与核查计划}

\section{辐射源设备监测}

\begin{table}[H]
\centering
\begin{tabular}{|c|c|c|c|}
\hline
监测类型 & 频次 & 执行单位 & 要求 \\
\hline
验收检测 & 新装 / 大修后 & 第三方资质机构 & 项目竣工控评验收检测 \\
\hline
状态检测 & 每年 1 次 & 第三方机构 & 年度设备性能检测 \\
\hline
稳定性自检 & 换源 / 维修后 & 院内物理师 & 自主质控检测 \\
\hline
\end{tabular}
\end{table}


\textbf{质控仪器配置}：放疗剂量仪、井型电离室、便携式 Xγ 辐射仪、个人剂量报警仪、出源标尺等，使用前送检校准。

\section{工作场所监测}

\begin{enumerate}
\item 第三方：每年 1 次全场所辐射水平检测；
\item 院内自主：物理师每月用便携式辐射仪巡检机房墙外各关注点，超标立即停机整改。
\end{enumerate}

\section{个人剂量监测}

\begin{table}[H]
\centering
\begin{tabular}{|c|c|}
\hline
项目 & 内容 \\
\hline
监测对象 & 全部 5 名放射工作人员 \\
\hline
监测周期 & 每 3 个月 1 次（合规≤3 个月） \\
\hline
受托单位 & 江西辐射剂量检测院有限公司 \\
\hline
管控目标 & 年均有效剂量＜5mSv 内控值 \\
\hline
\end{tabular}
\end{table}


\section{小结}

监测体系（第三方 + 院内自检）完善；建议全部质控设备按期检定、人员熟练掌握仪器使用、留存完整监测记录。
